% YY14200B
% Einschätzung der Indikation zur Wirbelsäulenimmobilisation
% 14.0
% 2013-11-30
\label{m:einschaetzung-indikation-ws-immobilisation}
\label{proc:YY14200B}
Eine Immobilisation der Wirbelsäule ist notwendig bei
Traumapatienten, bei denen folgendes zutrifft:

\begin{enumerate}
  \item \EE{Typischer Unfallmechanismus}
  \begin{compactitem}
    \item Hochgeschwindigkeitsunfall
    \item Sturz $\geq$ 3\Meter* oder 3fache Patientengröße
    \item Penetrierende Verletzungen im Bereich der
    Wirbelsäule
    \item Sportverletzungen im Kopf- oder Nackenbereich
    \item Trauma nach Sprung ins Wasser
    \item Stauchungsverletzung der HWS (z.\,B. Schlag auf
    Kopf)
    \item Suizidversuch durch Erhängen
    \item Bewusstloser Traumapatient
  \end{compactitem}
  \item \EE{Unklarer Unfallmechanismus}, bei dem einer der
  folgenden Punkte zutrifft:
  \begin{compactitem}
    \item Spinaler Schmerz oder Druckschmerz über der
    Wirbelsäule
    \item Auffälliger Befund bei Untersuchung von Motorik
    und Sensibilität\footnote{sofern nicht anders plausibel
    erklärbar}
    \item Patientenangaben kann nicht vertraut
    werden
    \begin{compactitem}
      \item Stresssituation
      \item Schädel- oder zerebrale Verletzung
      \item Auffälliger Bewusstseinszustand
      \item Intoxikation
      \item Ablenkende Verletzungen
    \end{compactitem}
  \end{compactitem}
  \item \EE{Verdacht}
\end{enumerate}

\noindent Solange es noch unklar ist, ob eine WS-Immobilisation
notwendig ist, muss eine manuelle Fixierung der HWS
durchgeführt werden!

\cite{ITLSde6,ScholzNotfallmedizin2,SiewertChirurgie8}\Commentary[Konsensus]{Review 2011-01-11}